\begin{remark*}[Exposition]
The preceding metrics become metric spaces once their domains are paired with
the verified distance functions.  This section records the resulting spaces as
objects in their own right.
\end{remark*}

\begin{definitionbox}{Definition --- Discrete Metric Space}
\begin{definition}[Discrete Metric Space]\label{def:discrete-metric-space}
Let \(X\) be a nonempty set.  The \emph{discrete metric space on \(X\)} is the
metric space
\[
  (X,d_{\mathrm{disc}}),
\]
where \(d_{\mathrm{disc}}\) is the discrete metric on \(X\).
\end{definition}
\end{definitionbox}

\begin{remark*}[Standard quantified statement]
\[
X\neq\varnothing
\Longrightarrow
\operatorname{MetricSpace}(X,d_{\mathrm{disc}}).
\]
\end{remark*}

\begin{remark*}[Interpretation]
A discrete metric space is a space where distinct points are all separated by
the same positive distance.  Its geometry records equality and inequality, but
no finer notion of nearness.
\end{remark*}

\begin{dependencies}
\begin{itemize}
  \item \hyperref[def:metric-space]{Definition: Metric Space}
  \item \hyperref[def:discrete-metric]{Definition: Discrete Metric}
  \item \hyperref[prop:discrete-metric-space]{Proposition: The Discrete Distance is a Metric}
\end{itemize}
\end{dependencies}

\begin{definitionbox}{Definition --- Euclidean Metric Space}
\begin{definition}[Euclidean Metric Space]\label{def:euclidean-metric-space}
For \(n\in\mathbb{N}\), the \emph{standard Euclidean metric space} is
\[
  (\mathbb{R}^n,d_2),
\]
where \(d_2\) is the Euclidean metric on \(\mathbb{R}^n\).
\end{definition}
\end{definitionbox}

\begin{remark*}[Standard quantified statement]
\[
n\in\mathbb{N}
\Longrightarrow
\operatorname{MetricSpace}(\mathbb{R}^n,d_2).
\]
\end{remark*}

\begin{remark*}[Interpretation]
The Euclidean metric space is the standard finite-dimensional distance model:
points are coordinate vectors, and distance is ordinary straight-line length.
\end{remark*}

\begin{dependencies}
\begin{itemize}
  \item \hyperref[def:metric-space]{Definition: Metric Space}
  \item \hyperref[def:euclidean-space]{Definition: Euclidean Space}
  \item \hyperref[def:euclidean-distance-rn]{Definition: Euclidean Distance on \(\mathbb{R}^n\)}
  \item \hyperref[cor:euclidean-metric]{Corollary: Euclidean Metric}
\end{itemize}
\end{dependencies}

\begin{definitionbox}{Definition --- Taxicab Metric Space}
\begin{definition}[Taxicab Metric Space]\label{def:taxicab-metric-space}
For \(n\in\mathbb{N}\), the \emph{taxicab metric space} is
\[
  (\mathbb{R}^n,d_1),
\]
where \(d_1\) is the taxicab metric on \(\mathbb{R}^n\).
\end{definition}
\end{definitionbox}

\begin{remark*}[Standard quantified statement]
\[
n\in\mathbb{N}
\Longrightarrow
\operatorname{MetricSpace}(\mathbb{R}^n,d_1).
\]
\end{remark*}

\begin{remark*}[Interpretation]
The taxicab metric space has the same underlying set as Euclidean space, but
distance is measured by summing coordinatewise displacements.
\end{remark*}

\begin{dependencies}
\begin{itemize}
  \item \hyperref[def:metric-space]{Definition: Metric Space}
  \item \hyperref[def:taxicab-metric-rn]{Definition: Taxicab Metric}
  \item \hyperref[prop:taxicab-metric-rn]{Proposition: Taxicab Metric}
\end{itemize}
\end{dependencies}

\begin{definitionbox}{Definition --- \(p\)-Metric Space}
\begin{definition}[\(p\)-Metric Space]\label{def:lp-metric-space}
Let \(1\leq p<\infty\) and let \(n\in\mathbb{N}\).  The \emph{\(p\)-metric
space} is
\[
  (\mathbb{R}^n,d_p),
\]
where \(d_p\) is the \(p\)-metric on \(\mathbb{R}^n\).
\end{definition}
\end{definitionbox}

\begin{remark*}[Standard quantified statement]
\[
1\leq p<\infty,\; n\in\mathbb{N}
\Longrightarrow
\operatorname{MetricSpace}(\mathbb{R}^n,d_p).
\]
\end{remark*}

\begin{remark*}[Interpretation]
The \(p\)-metric spaces form a family of finite-dimensional metric spaces on
the same coordinate set, with different rules for turning coordinate
displacements into distance.
\end{remark*}

\begin{dependencies}
\begin{itemize}
  \item \hyperref[def:metric-space]{Definition: Metric Space}
  \item \hyperref[def:lp-metric-rn]{Definition: \(p\)-Metric on \(\mathbb{R}^n\)}
  \item \hyperref[thm:lp-metric-rn]{Theorem: \(d_p\) is a Metric}
\end{itemize}
\end{dependencies}

\begin{definitionbox}{Definition --- Complex Metric Space}
\begin{definition}[Complex Metric Space]\label{def:complex-metric-space}
The \emph{complex metric space} is
\[
  (\mathbb{C},d_{\mathbb{C}}),
\]
where \(d_{\mathbb{C}}\) is the usual Euclidean distance on the complex plane.
\end{definition}
\end{definitionbox}

\begin{remark*}[Standard quantified statement]
\[
\operatorname{MetricSpace}(\mathbb{C},d_{\mathbb{C}}).
\]
\end{remark*}

\begin{remark*}[Interpretation]
The complex metric space treats complex numbers as points in the Euclidean
plane, while retaining the notation and algebra of \(\mathbb{C}\).
\end{remark*}

\begin{dependencies}
\begin{itemize}
  \item \hyperref[def:metric-space]{Definition: Metric Space}
  \item \hyperref[def:complex-metric]{Definition: Complex Metric}
  \item \hyperref[prop:complex-plane-metric-space]{Proposition: The Complex Plane is a Metric Space}
\end{itemize}
\end{dependencies}

\begin{definitionbox}{Definition --- Bounded Sequence Metric Space}
\begin{definition}[Bounded Sequence Metric Space]
\label{def:bounded-sequence-metric-space}
The \emph{bounded sequence metric space} is
\[
  (\ell^\infty(\mathbb{R}),d_\infty),
\]
where \(\ell^\infty(\mathbb{R})\) is the set of bounded real sequences and
\(d_\infty\) is the supremum metric.
\end{definition}
\end{definitionbox}

\begin{remark*}[Standard quantified statement]
\[
\operatorname{MetricSpace}(\ell^\infty(\mathbb{R}),d_\infty).
\]
\end{remark*}

\begin{remark*}[Interpretation]
The bounded sequence metric space is an infinite-dimensional example.  Its
distance measures the largest coordinatewise discrepancy between two bounded
real sequences.
\end{remark*}

\begin{dependencies}
\begin{itemize}
  \item \hyperref[def:metric-space]{Definition: Metric Space}
  \item \hyperref[def:bounded-real-sequence-space]{Definition: Supremum Metric on Bounded Real Sequences}
  \item \hyperref[prop:bounded-real-sequence-space-sup-metric]{Proposition: Supremum Metric on Bounded Sequences}
\end{itemize}
\end{dependencies}
